\chapter{Theoretische Grundlagen}
\label{chap:grundlagen}

\section{Eingebettete Systeme}
\label{sec:eingebettete-systeme}

Eingebettete Systeme werden von Marwedel als informationsverarbeitende Systeme 
definiert, die in in ein umgebendes Produkt eingebettet sind und dort dedizierte 
Steuerungs-, Regelungs- oder Überwachungsaufgaben übernehmen \cite{marwedel2021}. 
Im Gegensatz zu universellen Computersystemen sind sie auf einen konkreten 
Anwendungsfall zugeschnitten. Als gemeinsame Merkmale dieser Systemklasse nennt 
Marwedel insbesondere Echtzeitanforderungen sowie Anforderungen an Verlässlichkeit 
und Effizienz \cite{marwedel2021}, letztere äußert sich in ressourcensparender 
Hardware und häufig energiesparenden Betriebsweisen.

Die zunehmende Miniaturisierung und Leistungsfähigkeit moderner Mikrocontroller hat in 
den vergangenen Jahren zu einer breiten Durchdringung eingebetteter Systeme in 
unterschiedlichen Domänen geführt. Insbesondere die 
Integration drahtloser Kommunikationsschnittstellen wie WLAN und Bluetooth in 
kostengünstige Mikrocontroller-Plattformen eröffnet neue Möglichkeiten für verteilte, 
vernetzte Systemarchitekturen.

Im Kontext des entwickelten LED-Tube-Systems bilden eingebettete Systeme die zentrale 
Komponente sowohl der Bridge als auch der einzelnen Tubes. Beide setzen auf den 
ESP32-Mikrocontroller, der im Folgenden näher betrachtet wird.

\subsection{ESP32-Mikrocontroller}
\label{subsec:esp32}

Die Bezeichnung ESP32 steht nicht für einen einzelnen Chip, sondern für eine von 
Espressif Systems entwickelte Familie von System-on-Chips (SoC), deren Mitglieder sich 
in Prozessorarchitektur, Speicherausstattung und Peripherie unterscheiden 
\cite{espressif_esp32}. Kennzeichnend für die gesamte Familie ist die Integration eines 
leistungsfähigen Prozessors und umfangreicher drahtloser Kommunikationsschnittstellen 
(WLAN 802.11 b/g/n sowie Bluetooth beziehungsweise Bluetooth Low Energy) auf einem 
einzigen Baustein. Das namensgebende erste Mitglied der Familie kombiniert dabei einen 
Dual-Core-Prozessor des Typs Xtensa LX6 mit WLAN und Bluetooth 4.2. Seit der Einführung 
im Jahr 2016 hat sich die ESP32-Reihe zu einer der meistverwendeten Plattformen im 
Bereich des Internet of Things (IoT) und eingebetteter Systeme entwickelt. Besonders die 
Kompatibilität mit dem im Hobby- und Maker-Bereich weit verbreiteten Arduino-Ökosystem 
sowie die umfassende Unterstützung durch die Entwicklergemeinschaft haben zur 
Popularität des ESP32 beigetragen.

\subsubsection{Technische Merkmale}
\label{subsubsec:esp32-merkmale}

Es existieren verschiedene Varianten des ESP32, die sich in Bezug auf Speicherkapazität, Anzahl der GPIO-Pins,
integrierte Peripheriefunktionen und unterstützte Kommunikationsschnittstellen unterscheiden.
Die in diesem Projekt verwendeten Varianten verfügen über zwei unabhängige 32-Bit-Prozessorkerne mit einer maximalen 
Taktfrequenz von bis zu 240\,MHz. Diese Dual-Core-Architektur ermöglicht es, 
rechenintensive Aufgaben wie Signalverarbeitung und drahtlose Kommunikation parallel 
auszuführen.

Neben den Kommunikationsschnittstellen bieten die Bausteine der Familie umfangreiche 
Peripheriefunktionen, darunter mehrere UART-, SPI- und I²C-Schnittstellen, 
GPIO-Pins (General Purpose Input/Output), ADCs (Analog-Digital-Wandler), PWM-Kanäle 
(Pulse Width Modulation) sowie kapazitive Berührungssensoren \cite{espressif_esp32}. 

\subsubsection{Software-Entwicklung und Ökosystem}
\label{subsubsec:esp32-software}

Espressif Systems stellt mit dem ESP-IDF (Espressif IoT Development Framework) eine 
umfassende Entwicklungsumgebung bereit, die auf dem FreeRTOS-Echtzeitbetriebssystem 
basiert und native Unterstützung für C/C++-Programmierung bietet. Daneben existiert eine 
breite Community-Unterstützung für alternative Entwicklungsumgebungen wie Arduino IDE, 
PlatformIO sowie MicroPython.

Die Integration des ESP32 in Entwicklungsumgebungen sowie die Verfügbarkeit 
zahlreicher Open-Source-Bibliotheken haben die Einstiegshürde deutlich gesenkt und den 
Einsatz des Mikrocontrollers sowohl im Amateur- als auch im industriellen Kontext 
gefördert.

\subsubsection{Einsatz im entwickelten System}
\label{subsubsec:esp32-system}

Im Rahmen dieser Arbeit bildet der ESP32 die hardwaretechnische Grundlage sowohl für die 
zentrale Bridge als auch für die dezentralen LED-Tubes. Die Dual-Core-Architektur 
ermöglicht die parallele Verarbeitung von Netzwerkkommunikation (ESP-NOW, Art-Net) auf 
einem Kern und die Ansteuerung der LED-Hardware auf dem zweiten Kern.

Die integrierte WLAN-Funktionalität wird für die Implementierung des proprietären 
ESP-NOW-Protokolls genutzt, welches die drahtlose Datenübertragung zwischen Bridge und 
Tubes realisiert. Gleichzeitig erlaubt die Ethernet-Erweiterbarkeit über externe Module 
die Anbindung der Bridge an professionelle Lichtsteuersysteme mittels Art-Net.

Die GPIO-Pins dienen der Ansteuerung der Peripheriekomponenten, etwa der LED-Strips und 
der übrigen in der Schaltung verbauten ICs (Integrated Circuits).
Die Leistungsfähigkeit des ESP32 stellt sicher, dass die zeitkritischen Anforderungen der Lichtsteuerung erfüllt werden können, 
während gleichzeitig die Flexibilität für eventuelle zukünftige Erweiterungen gegeben ist.

\section{Kommunikationsprotokolle}
\label{sec:kommunikationsprotokolle}

\subsection{DMX512}
\label{subsec:dmx512}

DMX512 (Digital Multiplex 512) ist ein weit verbreiteter Industriestandard für die
Steuerung von Bühnen- und Veranstaltungstechnik. Der Standard wurde 1986 vom United 
States Institute for Theatre Technology (USITT) entwickelt und wird heute von der 
Entertainment Services and Technology Association (ESTA) gepflegt. Er wurde seither in 
mehreren Revisionen fortgeschrieben die aktuelle Fassung ist ANSI E1.11-2024 
\cite{ANSI_E1.11_2024}.

\subsubsection{Technische Grundlagen}
\label{subsubsec:dmx512-technik}

DMX512 überträgt Daten seriell und unidirektional über eine differentielle
RS-485-Leitung mit einer Datenrate von 250\,kbit/s \cite{ANSI_E1.11_2024}. Die
Übertragung erfolgt in Form von Paketen, die neben dem einleitenden START-Code bis zu
512 Datenfelder (Slots) umfassen. Jeder Slot wird durch einen 8-Bit-Wert (0--255)
repräsentiert, der beispielsweise die Helligkeit, Farbe oder Position eines Geräts
definiert.

Die physikalische Verbindung erfolgt über eine Daisy-Chain-Topologie, bei der die Geräte
in Reihe geschaltet werden. Der Standard schreibt vor, die Leitung abzuschließen, um
Reflexionen und die daraus folgenden Fehlfunktionen zu vermeiden als Abschluss ist ein
Widerstand von 120\,Ohm vorgesehen. Die Belastung einer DMX512-Linie ist zudem auf 32
Unit Loads begrenzt, wobei jeder Empfängereingang höchstens eine Unit Load darstellen
darf \cite{ANSI_E1.11_2024}. Größere Installationen erfordern daher Verstärker
(DMX-Splitter), die zugleich die nutzbare Leitungslänge erhöhen. Eine maximale
Leitungslänge legt der Standard selbst nicht fest, sie ergibt sich aus den elektrischen
Eigenschaften der RS-485-Strecke und liegt in der Praxis im Bereich einiger hundert
Meter.

\subsubsection{Protokollstruktur}
\label{subsubsec:dmx512-protokoll}

Ein DMX512-Paket beginnt mit einer Reset-Sequenz aus \textit{Break}, einer Unterbrechung
der Datenleitung, dem darauf folgenden \textit{Mark After Break} (MAB) und dem
\textit{START-Code}. Der START-Code kennzeichnet die Art der nachfolgenden Daten für
Beleuchtungsdaten ist der NULL-START-Code mit dem Wert 0x00 vorgesehen. Anschließend
folgen bis zu 512 Datenslots. Für die Zeiten der Reset-Sequenz unterscheidet der
Standard zwischen Sende- und Empfangsseite: Sender müssen einen Break von mindestens
92\,µs und einen MAB von mindestens 12\,µs erzeugen, während Empfänger bereits einen
Break ab 88\,µs und einen MAB ab 8\,µs als Paketbeginn erkennen müssen
\cite{ANSI_E1.11_2024}. Für die Implementierung des Empfängers ist somit die geringere,
für die des Senders die höhere Vorgabe maßgeblich.

Aus Datenrate und Paketaufbau ergibt sich für ein vollständiges Paket mit 513 Slots eine
minimale Aktualisierungszeit von 22,7\,ms und damit eine maximale Wiederholrate von
44\,Hz \cite{ANSI_E1.11_2024} werden weniger Slots übertragen, steigt die erreichbare
Rate entsprechend. Eine bidirektionale Kommunikation ist im regulären Datenpfad nicht
vorgesehen, sodass der Sender keine Rückmeldung über den Status der Empfangsgeräte
erhält.

\subsubsection{Einsatz in der Veranstaltungstechnik}
\label{subsubsec:dmx512-einsatz}

DMX512 hat sich als De-facto-Standard in der professionellen Lichtsteuerung etabliert
und wird von nahezu allen Herstellern unterstützt. Die Einfachheit des Protokolls sowie
die deterministischen Übertragungseigenschaften machen es besonders geeignet für
zeitkritische Anwendungen, bei denen eine zuverlässige Steuerung
erforderlich ist.

Im Kontext des entwickelten LED-Tube-Systems dient der Aufbau des DMX512-Pakets als
Vorbild für das interne Datenformat, in dem die Steuerbefehle repräsentiert werden. Die Bridge kann DMX-Daten sowohl über Art-Net 
(Ethernet) als auch über eine DMX-Schnittstelle (XLR) empfangen und verarbeitet 
diese anschließend zur Übertragung an die LED-Tubes. Dabei bildet die Struktur des 
DMX512-Frames die Grundlage für das entwickelte drahtlose Übertragungsprotokoll, das 
über ESP-NOW an die Tubes übertragen wird. 
Die Orientierung an diesem etablierten Standard ermöglicht eine nahtlose Integration in 
bestehende Lichtsteuersysteme und gewährleistet Kompatibilität mit professioneller 
Steuerungssoftware.


\subsection{Art-Net}
\label{subsec:artnet}

Art-Net ist ein Ethernet-basiertes Kommunikationsprotokoll, das von Artistic Licence
entwickelt wurde und die Übertragung von DMX512-Daten über IP-Netzwerke ermöglicht.
Seit seiner Einführung im Jahr 1998 hat sich Art-Net zu einem weit verbreiteten Standard
in der professionellen Lichtsteuerung entwickelt und wird von zahlreichen
Softwarelösungen und Hardwarekomponenten unterstützt.

Die Etablierung von Art-Net als De-facto-Standard in der professionellen 
Veranstaltungstechnik ist auf die Limitierung klassischer DMX512-Verbindungen 
zurückzuführen. Während eine einzelne DMX-Linie lediglich 512 Kanäle/Byte bereitstellt, 
erfordern moderne Installationen mit einer Vielzahl komplexer Scheinwerfer, 
LED-Pixelmatrizen oder hochauflösenden LED-Wänden häufig mehrere tausend Steuerkanäle. 
Die Art-Net-4-Spezifikation sieht ein theoretisches Maximum von 32.768 DMX-Universen 
vor, wie viele davon tatsächlich übertragen werden können, hängt von der verfügbaren 
Netzwerkbandbreite und der Leistungsfähigkeit der beteiligten Geräte ab 
\cite{artnet4}. Art-Net adressiert damit die gestiegenen Bandbreitenanforderungen 
professioneller Lichtinstallationen. Insbesondere bei der Ansteuerung von LED-Pixelmonitoren, 
bei denen jeder einzelne Pixel mehrere DMX-Kanäle (typischerweise RGB oder RGBW) 
beansprucht, stößt klassisches DMX512 schnell an seine Grenzen, während Art-Net eine 
skalierbare Lösung bietet.

\subsubsection{Funktionsweise und Architektur}
\label{subsubsec:artnet-funktion}

Art-Net kapselt DMX512-Daten in UDP-Pakete (User Datagram Protocol) und überträgt diese
über ein Standard-Ethernet-Netzwerk. Während eine DMX512-Linie auf ein einzelnes
RS-485-Segment mit einem Universum von 512 Kanälen beschränkt ist, ermöglicht Art-Net die
gleichzeitige Übertragung mehrerer Universen über eine einzige Netzwerkverbindung.

Ein Art-Net-Netzwerk besteht typischerweise aus einem oder mehreren \textit{Controllern}
(z.\,B. Lichtsteuerungssoftware), die Art-Net-Pakete senden, sowie \textit{Nodes}, die
diese Pakete empfangen und in DMX512-Signale umwandeln. Die Nodes können dabei sowohl
als Sender als auch als Empfänger fungieren, wodurch bidirektionale Kommunikation
möglich ist.

\subsubsection{Protokollstruktur}
\label{subsubsec:artnet-protokoll}

Jedes Art-Net-Paket beginnt mit einer acht Byte langen, nullterminierten Kennung, 
gefolgt von einem OpCode, der den Pakettyp bestimmt, der
Protokollversion sowie weiteren Verwaltungsinformationen \cite{artnet4}. Das Ziel wird
über die Port-Address adressiert, eine 15\,Bit breite Zahl, die sich aus den
Feldern Net, Sub-Net und Universe34 zusammensetzt und die 32.768 möglichen Universen
abbildet \cite{artnet4}. Der Nutzdatenbereich umfasst bis zu 512 Byte, entsprechend
einem DMX512-Universum.

\subsubsection{Vorteile gegenüber klassischem DMX512}
\label{subsubsec:artnet-vorteile}

Durch die Nutzung von Ethernet-Infrastruktur bietet Art-Net mehrere Vorteile gegenüber
klassischem DMX512:
\begin{itemize}
    \item \textbf{Skalierbarkeit:} Mehrere DMX-Universen können parallel über eine
          einzige Netzwerkverbindung übertragen werden.
    \item \textbf{Reichweite:} Die Ausdehnung ist nicht mehr durch die Grenzen eines
          einzelnen RS-485-Segments bestimmt, sondern kann mit
          Standard-Netzwerkkomponenten nahezu beliebig erweitert werden.
    \item \textbf{Bidirektionalität:} Rückmeldungen über Gerätestatus und Fehlerzustände sind möglich.
    \item \textbf{Flexibilität:} Bestehende IT-Infrastrukturen können genutzt werden,
          wodurch separate DMX-Verkabelung entfällt.
\end{itemize}

\subsubsection{Einsatz im entwickelten System}
\label{subsubsec:artnet-system}

Im entwickelten LED-Tube-System dient Art-Net als Schnittstelle zur Anbindung an
professionelle Lichtsteuerungssoftware. Die Bridge empfängt Art-Net-Pakete über ihre
Ethernet-Schnittstelle und extrahiert daraus die DMX512-Daten, die anschließend wie im
kabelgebundenen DMX-Betrieb weiterverarbeitet werden.

Darüber hinaus ist ein kabelgebundener Pixelsteuerungsmodus vorgesehen, in dem die
einzelnen Pixel der LED-Strips über eine Datenleitung angesprochen werden. Die Hardware
der Bridge ist dafür bereits vorbereitet
(Abschnitt~\ref{subsec:bridge-pcb}), die Umsetzung dieses Modus ist jedoch nicht
Gegenstand der vorliegenden Arbeit. Durch die Unterstützung von Art-Net bleibt das System
mit einer Vielzahl bestehender Steuerungslösungen kompatibel.

\section{Drahtlose Kommunikation}
\label{sec:drahtlose-kommunikation}

Drahtlose Kommunikation stellt einen zentralen Bestandteil moderner eingebetteter
Systeme dar, insbesondere im Kontext verteilter Architekturen und modularer
Installationen. Im Vergleich zu kabelgebundenen Kommunikationsformen bietet sie eine
höhere Flexibilität und vereinfacht die physische Integration einzelner Komponenten.
Gleichzeitig ergeben sich jedoch neue Herausforderungen hinsichtlich Latenz,
Zuverlässigkeit, Störanfälligkeit und Energieeffizienz.

Im vorliegenden System wird drahtlose Kommunikation eingesetzt, um Steuerdaten von einer
zentralen Bridge an mehrere räumlich verteilte LED-Tubes zu übertragen. Zunächst werden
das genutzte Frequenzband und der Zugriff auf das Übertragungsmedium betrachtet, da
beide das Verhalten der Funkstrecke unter Last bestimmen. anschließend wird das
eingesetzte Protokoll ESP-NOW vorgestellt.

\subsection{Das 2,4-GHz-Band und der Zugriff auf das Übertragungsmedium}
\label{subsec:24ghz-medienzugriff}

Die Übertragung erfolgt im Frequenzbereich von 2400 bis 2483,5\,MHz. Dieses Band steht
weltweit anmeldefrei zur Verfügung. in Europa regelt die harmonisierte Norm
ETSI EN 300 328 die Nutzung durch Breitband-Datenübertragungssysteme
\cite{etsi_en300328}. Genutzt wird es nicht exklusiv, sondern gemeinsam von WLAN,
Bluetooth, Zigbee und weiteren Funkanwendungen, die die genannte Norm ausdrücklich
gemeinsam erfasst. Für ein Steuerungssystem
folgt daraus, dass die verfügbare Übertragungskapazität nicht allein von der eigenen
Hardware abhängt, sondern von der Belegung des Kanals durch fremde Sender.

Innerhalb des Bands definiert IEEE 802.11 in Europa dreizehn Kanäle im Abstand von je
5\,MHz, während ein Kanal eine Bandbreite von etwa 20\,MHz belegt. 

Der Zugriff auf das gemeinsam genutzte Medium erfolgt nach dem in IEEE 802.11
festgelegten Verfahren CSMA/CA (Carrier Sense Multiple Access with Collision
Avoidance) \cite{ieee80211_2024}. Eine Station prüft vor dem Senden, ob der Kanal frei
ist, wartet bei belegtem Kanal ab und beginnt die Übertragung erst nach einer zusätzlich
zufällig gewählten Wartezeit (Backoff). Kollisionen werden dadurch vermieden, gleichzeitig verzögern sich
Sendevorgänge bei hoher Fremdlast oder unterbleiben ganz. Für die Bewertung der
Zuverlässigkeit ist diese Unterscheidung wesentlich: Eine ausgelastete Umgebung führt
nicht zu verfälschten, sondern zu verspäteten oder ausbleibenden Paketen. Die
messtechnische Untersuchung dieses Effekts erfolgt in
Abschnitt~\ref{sec:stoerfestigkeit}.

\subsection{ESP-NOW}
\label{subsec:espnow-protokoll}
Das von Espressif Systems entwickelte ESP-NOW-Protokoll stellt eine proprietäre, verbindungslose 
Kommunikationstechnologie dar, die primär für Low-Power-Anwendungen im Bereich des Internet of Things (IoT) konzipiert wurde. 
Im Gegensatz zu herkömmlichen WLAN-Verbindungen nach dem IEEE 802.11-Standard, welche eine zeit- und 
energieintensive Assoziierung mit einem Access Point sowie den vollständigen Durchlauf des TCP/IP-Stacks erfordern, 
operiert ESP-NOW unmittelbar auf der Sicherungsschicht (Data Link Layer) des ISO-OSI-Modells.

Technisch basiert die Übertragung auf sogenannten \textit{IEEE 802.11 Vendor Specific Action Frames} \cite{espressif_espnow}. 
Diese ermöglichen es den Knotenpunkten, Datenpakete direkt über die MAC-Adresse zu adressieren, ohne dass ein 
vorheriger Handshake oder eine IP-Zuweisung notwendig ist. Daraus resultiert eine signifikante Reduktion 
der Latenzzeit sowie des energetischen Overheads, da die Funkmodule lediglich für die Dauer der Übertragung aktiviert werden müssen.

Trotz der verbindungslosen Natur unterstützt das Protokoll eine optionale bidirektionale 
Kommunikation durch Bestätigungs-Frames (Acknowledgements) sowie eine Verschlüsselung der 
Action Frames \cite{espressif_espnow}. Die maximale Nutzlast 
je Datenpaket beträgt in der ersten Version 250 Byte und wurde in der zweiten 
Version auf 1470 Byte erhöht \cite{espressif_espnow}. Mit Version 1.0 lässt sich somit nur 
etwa ein halber DMX512-Frame übertragen, während Version 2.0 ein vollständiges 
DMX-Universum in einem einzigen Paket aufnehmen kann.

Version 2 ist abwärtskompatibel zu Version 1, sodass ältere Geräte weiterhin mit der neuen Version kommunizieren können, 
jedoch nur die geringere Payload-Größe von 250 Byte nutzen können.
Die Standard-Übertragungsrate gibt Espressif mit 1\,Mbit/s an \cite{espressif_espnow}.

ESP-NOW verwendet ein Peer-to-Peer- beziehungsweise Broadcast-Kommunikationsmodell, bei
dem Datenpakete direkt zwischen den beteiligten Knoten ausgetauscht werden. Für die
gerichtete Übertragung muss jeder Empfänger zuvor als Peer registriert werden, wobei die
Zahl gleichzeitig registrierter Peers auf 20 begrenzt ist \cite{espressif_espnow}. Diese
Begrenzung entfällt im Broadcast-Betrieb, da dort sämtliche Empfänger über eine einzige
Broadcast-Adresse erreicht werden.

\subsubsection{Latenz, Zuverlässigkeit und Einschränkungen}
\label{subsubsec:espnow-eigenschaften}

Durch den Verzicht auf den vollständigen WLAN-Stack weist ESP-NOW in der Regel eine
geringere Latenz als klassische WLAN-Kommunikation auf. Gleichzeitig existieren jedoch
keine integrierten Mechanismen zur garantierten Zustellung oder zur Reihenfolgekontrolle
von Paketen. Der Sendestatus, den die Anwendung nach einem Sendevorgang erhält, bezieht
sich ausschließlich auf die MAC-Schicht, ob die Daten die Anwendungsschicht des
Empfängers erreichen, ist damit ausdrücklich nicht sichergestellt
\cite{espressif_espnow}.

Die Zuverlässigkeit der Übertragung hängt daher stark von den Umgebungsbedingungen, der
Anzahl der beteiligten Knoten sowie der Paketfrequenz ab. Da ESP-NOW auf 802.11-Frames
aufsetzt, unterliegt es zudem dem in
Abschnitt~\ref{subsec:24ghz-medienzugriff} beschriebenen Medienzugriff: Bei hoher
Kanalauslastung verzögern sich Sendevorgänge oder unterbleiben ganz. Paketverluste sind
somit prinzipiell möglich und müssen auf Anwendungsebene berücksichtigt werden.

\subsubsection{Bedeutung für das entwickelte System}
\label{subsubsec:espnow-system}

ESP-NOW bildet die Grundlage für die drahtlose Übertragung von DMX-Steuerdaten von der
Bridge zu den LED-Tubes. Die geringe Latenz und der geringe Overhead machen das
Protokoll besonders geeignet für zeitkritische Steuerinformationen, wie sie in der
Lichtsteuerung erforderlich sind.

Gleichzeitig stellt die fehlende Zustellgarantie eine zentrale Herausforderung dar, die
im Rahmen dieser Arbeit durch experimentelle Untersuchungen der Verbindungsstabilität
analysiert wird. Die Ergebnisse dieser Untersuchungen bilden die Basis für die Bewertung
der Eignung von ESP-NOW als drahtlose Alternative zu kabelgebundenen DMX-Verbindungen.
